\begin{pastexam}[label={px:ntu-109-comm-dsa-03}]{臺大電信所 109 學年度（B）第 3 題}
（是非；每題 2 分，答錯倒扣 1 分）

In a weighted undirected graph $G=(V,E,w)$, breadth-first search from a vertex $s$ finds
single-source shortest paths from $s$ (via parent pointers) in $O(V+E)$ time.
\end{pastexam}

\begin{pastsoln}
答案：F（錯）。

BFS 完全不看權重，它找到的是\textbf{邊數}最少的路徑（資結書定理~21.4 的 $\delta$ 是邊數）。
有權重時，邊數少不代表權重和小。

反例：三個頂點 $s, a, t$，$w(s, a) = 1$、$w(a, t) = 1$、$w(s, t) = 5$。 從 $s$ 做 BFS，$t$
是 $s$ 的鄰居，$\pi[t] = s$，父指標給出的路徑是 $s, t$，權重 $5$；但 $s, a, t$ 的權重只有
$2$。

有權重的單源最短路徑要用 Dijkstra（權重非負）或 Bellman--Ford（資結書第~25.2、25.3~節）。
只有所有權重都相同時，BFS 才剛好是答案。
\end{pastsoln}
